% GENERATED FILE. Edit canonical lessons, then run npm run generate:books.
% canonical-source-hash: fnv1a64:68126ade7abaa9ad
% canonical-lessons: IT-C373-purche, IT-C373-qualora, IT-C373-altrimenti, IT-C373-senno, IT-C373-eventuale

\chapter{A una condizione}
\label{ch:it-a-una-condizione}
\hlchaptermodality{373}

\begin{chapteropening}
\textbf{By the end of this chapter:} \emph{I can set a condition: provided that, in case, otherwise, or else, and a possible extra cost.}

Complete the last lesson of chapter 373: I can set a condition: provided that, in case, otherwise, or else, and a possible extra cost.
\end{chapteropening}

\section[purché]{purché — provided that, as long as}

\label{lesson:IT-C373-purche}



\begin{quote}
Before the new one: say the Italian for plausible, then the Italian for the precondition.
\end{quote}

\subsection*{You'll want to know: purché}
\textbf{purché} — \textquotedblleft{}provided that, as long as\textquotedblright{}.

\begin{quote}
Puoi uscire, purché tu sia a casa prima di mezzanotte.
\end{quote}

\begin{quote}
\emph{You can go out, as long as you are home before midnight.}
\end{quote}

Always followed by the subjunctive: \emph{purché tu sia}, \emph{purché lui venga}.

\begin{cousinweb}
\emph{purché} is built from \emph{pure} \textquotedblleft{}simply, only\textquotedblright{}, from Latin \emph{pūrē} \textquotedblleft{}purely\textquotedblright{} (Lewis \& Short), on \emph{pūrus} \textquotedblleft{}clean, pure\textquotedblright{}, and \emph{che} \textquotedblleft{}that\textquotedblright{}: \textquotedblleft{}only that\textquotedblright{}.
\end{cousinweb}

\begin{grammarlens}[title={What you've built}]
One more word for telling, explaining and describing at B1.
\end{grammarlens}

\subsection*{Your turn}
\begin{itemize}
\raggedright
  \item \emph{Say it:} \emph{purché}
  \item \emph{Say it:} \emph{purché}, once more
  \item \emph{Say it:} a sentence of your own with \emph{purché}
  \item \emph{Recall:} say the Italian for plausible, then the Italian for the precondition, then say \emph{purché} again
\end{itemize}

\begin{tcolorbox}[breakable,colback=teal!4,colframe=teal!35!black,title={Before you move on}]
What does purché mean? (\textquotedblleft{}Provided that, as long as\textquotedblright{}.) Say it once more.
\end{tcolorbox}
\section[qualora]{qualora — if, in case (formal)}

\label{lesson:IT-C373-qualora}



\begin{quote}
Before the new one: say the Italian for the precondition, then the Italian for provided that.
\end{quote}

\subsection*{You'll want to know: qualora}
\textbf{qualora} — \textquotedblleft{}if, in case (formal)\textquotedblright{}.

\begin{quote}
Qualora non foste soddisfatti, vi restituiamo i soldi.
\end{quote}

\begin{quote}
\emph{Should you not be satisfied, we will give you your money back.}
\end{quote}

Formal and written, followed by the subjunctive. In speech say \emph{se}.

\begin{cousinweb}
\emph{qualora} is built from \emph{quale} \textquotedblleft{}which\textquotedblright{}, from Latin \emph{quālis}, and \emph{ora} \textquotedblleft{}hour\textquotedblright{}, from Latin \emph{hōra} (Lewis \& Short): \textquotedblleft{}at whatever hour\textquotedblright{}, and so \textquotedblleft{}if ever\textquotedblright{}.
\end{cousinweb}

\begin{grammarlens}[title={What you've built}]
One more word for telling, explaining and describing at B1.
\end{grammarlens}

\subsection*{Your turn}
\begin{itemize}
\raggedright
  \item \emph{Say it:} \emph{qualora}
  \item \emph{Say it:} \emph{qualora}, once more
  \item \emph{Say it:} a sentence of your own with \emph{qualora}
  \item \emph{Recall:} say the Italian for the precondition, then the Italian for provided that, then say \emph{qualora} again
\end{itemize}

\begin{tcolorbox}[breakable,colback=teal!4,colframe=teal!35!black,title={Before you move on}]
What does qualora mean? (\textquotedblleft{}If, in case (formal)\textquotedblright{}.) Say it once more.
\end{tcolorbox}
\section[altrimenti]{altrimenti — otherwise, or else}

\label{lesson:IT-C373-altrimenti}



\begin{quote}
Before the new one: say the Italian for provided that, then the Italian for if.
\end{quote}

\subsection*{You'll want to know: altrimenti}
\textbf{altrimenti} — \textquotedblleft{}otherwise, or else\textquotedblright{}.

\begin{quote}
Sbrigati, altrimenti perdiamo l'autobus.
\end{quote}

\begin{quote}
\emph{Hurry up, otherwise we will miss the bus.}
\end{quote}

Neutral in speech and writing.

\begin{cousinweb}
\emph{altrimenti} is built from \emph{altro} \textquotedblleft{}other\textquotedblright{} and the adverb ending \emph{-mente} (EtymDB); \emph{altro} is inherited from Latin \emph{alter} \textquotedblleft{}the other\textquotedblright{}, and \emph{-mente} from \emph{mēns, mentis} \textquotedblleft{}mind\textquotedblright{}: \textquotedblleft{}in another way\textquotedblright{}.
\end{cousinweb}

\begin{grammarlens}[title={What you've built}]
One more word for telling, explaining and describing at B1.
\end{grammarlens}

\subsection*{Your turn}
\begin{itemize}
\raggedright
  \item \emph{Say it:} \emph{altrimenti}
  \item \emph{Say it:} \emph{altrimenti}, once more
  \item \emph{Say it:} a sentence of your own with \emph{altrimenti}
  \item \emph{Recall:} say the Italian for provided that, then the Italian for if, then say \emph{altrimenti} again
\end{itemize}

\begin{tcolorbox}[breakable,colback=teal!4,colframe=teal!35!black,title={Before you move on}]
What does altrimenti mean? (\textquotedblleft{}Otherwise, or else\textquotedblright{}.) Say it once more.
\end{tcolorbox}
\section[sennò]{sennò — or else, otherwise (casual)}

\label{lesson:IT-C373-senno}



\begin{quote}
Before the new one: say the Italian for if, then the Italian for otherwise.
\end{quote}

\subsection*{You'll want to know: sennò}
\textbf{sennò} — \textquotedblleft{}or else, otherwise (casual)\textquotedblright{}.

\begin{quote}
Mettiti la giacca, sennò hai freddo.
\end{quote}

\begin{quote}
\emph{Put your jacket on, or else you will be cold.}
\end{quote}

Spoken and informal; also written \emph{se no}.

\begin{cousinweb}
\emph{sennò} is \emph{se no} \textquotedblleft{}if not\textquotedblright{} written as one word: \emph{se} is inherited from Latin \emph{sī} \textquotedblleft{}if\textquotedblright{} and \emph{no} from Latin \emph{nōn} \textquotedblleft{}not\textquotedblright{} (Lewis \& Short).
\end{cousinweb}

\begin{grammarlens}[title={What you've built}]
One more word for telling, explaining and describing at B1.
\end{grammarlens}

\subsection*{Your turn}
\begin{itemize}
\raggedright
  \item \emph{Say it:} \emph{sennò}
  \item \emph{Say it:} \emph{sennò}, once more
  \item \emph{Say it:} a sentence of your own with \emph{sennò}
  \item \emph{Recall:} say the Italian for if, then the Italian for otherwise, then say \emph{sennò} again
\end{itemize}

\begin{tcolorbox}[breakable,colback=teal!4,colframe=teal!35!black,title={Before you move on}]
What does sennò mean? (\textquotedblleft{}Or else, otherwise (casual)\textquotedblright{}.) Say it once more.
\end{tcolorbox}
\section[eventuale]{eventuale — possible, any (that may arise)}

\label{lesson:IT-C373-eventuale}



\begin{quote}
Before the new one: say the Italian for otherwise, then the Italian for or else.
\end{quote}

\subsection*{You'll want to know: eventuale}
\textbf{eventuale} — \textquotedblleft{}possible, any (that may arise)\textquotedblright{}.

\begin{quote}
Per eventuali domande, scrivete alla segreteria.
\end{quote}

\begin{quote}
\emph{For any questions, write to the office.}
\end{quote}

It means possible, not final: \emph{eventualmente} is if need be, not eventually.

\begin{cousinweb}
\emph{eventuale} is usually traced to French \emph{éventuel}, built on Latin \emph{ēventus} \textquotedblleft{}an outcome, an event\textquotedblright{} (Lewis \& Short), from \emph{ēvenīre} \textquotedblleft{}to come out, to happen\textquotedblright{}, on \emph{venīre} \textquotedblleft{}to come\textquotedblright{}.
\end{cousinweb}

\begin{grammarlens}[title={What you've built}]
One more word for telling, explaining and describing at B1.
\end{grammarlens}

\subsection*{Your turn}
\begin{itemize}
\raggedright
  \item \emph{Say it:} \emph{eventuale}
  \item \emph{Say it:} \emph{eventuale}, once more
  \item \emph{Say it:} a sentence of your own with \emph{eventuale}
  \item \emph{Recall:} say the Italian for otherwise, then the Italian for or else, then say \emph{eventuale} again
\end{itemize}

\begin{tcolorbox}[breakable,colback=teal!4,colframe=teal!35!black,title={Before you move on}]
What does eventuale mean? (\textquotedblleft{}Possible, any (that may arise)\textquotedblright{}.) Say it once more.
\end{tcolorbox}
